%! Elimination Matrices and Inverse Matrices — Unit 2, Lesson 2.2 — Warm-Up
\documentclass[10pt]{article}
\usepackage{linalg-article}
\usepackage{linalg-boxes}
\newcommand{\TallMath}[1]{$\displaystyle #1\rule[-1.4em]{0pt}{3.2em}$}
\newcommand{\xx}{\mathbf{x}}
\newcommand{\bb}{\mathbf{b}}

\begin{document}

\pageheader{Unit 2, Lesson 2.2}{Warm-Up}

\vspace{0.12in}

\begin{notesbox}{Getting Ready: Skills We'll Use Today}
\small Today every elimination step becomes a \emph{matrix}. These three moves are the pieces.
Answer quickly.

\vspace{4pt}
\textbf{1. Matrix times a vector} (Lesson 1.3). Multiply row by row:
\[
  \begin{bmatrix} 1 & 0 \\ -2 & 1 \end{bmatrix}
  \begin{bmatrix} 5 \\ 12 \end{bmatrix}
  = \begin{bmatrix}\rule{0pt}{1.1em}\blank{1.1cm}\\[2pt]\blank{1.1cm}\end{bmatrix}.
\]
\begin{work}
  1(5)+0(12) &= 5 \\
  -2(5)+1(12) &= 2
\end{work}

\vspace{4pt}
\textbf{2. Matrix times a matrix} (Lesson 1.4). Take the product one column at a time:
\[
  \begin{bmatrix} 1 & 0 \\ -3 & 1 \end{bmatrix}
  \begin{bmatrix} 2 & 1 \\ 6 & 5 \end{bmatrix}
  = \begin{bmatrix}\rule{0pt}{1.1em}\blank{1.1cm} & \blank{1.1cm}\\[2pt]\blank{1.1cm} & \blank{1.1cm}\end{bmatrix}.
\]
Notice the top row is unchanged, and the bottom row became (row~2)$\,-3\,$(row~1).
\begin{work}
  -3(2,1)+(6,5) &= (0,2)
\end{work}

\vspace{4pt}
\textbf{3. One multiplier} (Lesson 2.1). To clear the $x$ below the pivot in
$\ x + 3y = 7,\ \ 2x + 7y = 15$, the pivot is $1$ and the multiplier is
\[
  \ell = \frac{\text{entry to eliminate}}{\text{pivot}} = \frac{2}{1} = \blank{1.2cm}.
\]
\end{notesbox}

\end{document}
